\section{A robust tag clock}\label{clock:section}

The shift table supplies compatible data when both factors use the same
tag. The next construction limits the damage from independent tag
choices. For each possible total time, it will permit only equal tags
or exactly one unequal ordered pair. Later, the episode's affine shift
will accommodate that one exceptional pair.

We need more than this compatibility rule. A finite roster of candidate
positions must encounter few exceptional totals, uniformly under every
translation. Outside the exceptional set, each tag must admit a
decomposition with a common positive interior margin, so a sufficiently
fine grid can find one. The logarithmic estimate and its
forest argument are developed in \cite[Lemma~3.1]{HeightThirteen}, with
constant $1{,}300{,}000$. Here the estimate $10^7\log(2m)$ accompanies a
torus construction whose common interior margin and perturbation
tolerance let us postpone the final clock period.

Write $\mathbb T=\mathbb R/\mathbb Z$ and
$\|x\|_{\mathbb T}=\min_{j\in\mathbb Z}|x-j|$. A finite product
$\mathcal T=\mathbb T^C$ is given the maximum metric
$\operatorname{dist}(u,v)=\max_{c\in C}\|u_c-v_c\|_{\mathbb T}$.
Addition and subtraction are coordinatewise. For subsets, $I+J$ denotes
the set of sums of their elements, and $z-J=\{z-j:j\in J\}$.
All logarithms in the clock estimate are natural.

\begin{lemma}[Robust tag clock]\label{clock:robust}
Let $\Lambda$ and $\mathcal P$ be finite nonempty sets, with
$m=|\Lambda|$. There exist a finite-dimensional torus $\mathcal T$,
sets $I_\alpha,J_\alpha\subseteq\mathcal T$ for $\alpha\in\Lambda$,
a set $W\subseteq\mathcal T$, points $v_p\in\mathcal T$ for
$p\in\mathcal P$, and constants $\rho,\eta>0$ with these properties:
\begin{enumerate}[label=\textup{(\roman*)}]
\item For each $z\in\mathcal T$, the edge set
\begin{equation}\label{clock:edge-set}
 \mathcal E_z=\{(\alpha,\lambda)\in\Lambda^2:
                                  z\in I_\alpha+J_\lambda\}
\end{equation}
is either contained in the diagonal or consists of one off-diagonal pair.
\item If $z\notin W$, then for every $\alpha\in\Lambda$,
$I_\alpha\cap(z-J_\alpha)$ contains a ball of radius $\rho$.
In particular $\mathcal E_z$ is the full diagonal.
\item For every family $(\tilde v_p)_{p\in\mathcal P}$ satisfying
$\operatorname{dist}(\tilde v_p,v_p)<\eta$ for all $p$, and every
$a\in\mathcal T$,
\begin{equation}\label{clock:roster-estimate}
 \#\{(p,p')\in\mathcal P^2:p\ne p',\quad
               a+\tilde v_{p'}-\tilde v_p\in W\}
                         \le c_{\rm cl}\log(2m),
 \qquad c_{\rm cl}=10^7.
\end{equation}
\end{enumerate}
The constants $\rho,\eta$ and the dimension may depend on both finite
sets, but $c_{\rm cl}$ is absolute.
\end{lemma}

\begin{proof}
The proof separates compatibility from counting. First coordinates will
exclude all competitors of an unequal edge except its reversal. Type
batches will exclude the reversal too. Each block also permits a diagonal
decomposition away from an explicitly specified obstruction. Finally we
assign the roster values so that, at any translation, very few ordered
pairs can meet the union of those obstructions.

\paragraph{Partitions of the tags.}
Put
\[
        a_0=\lceil64\log(2m)\rceil,\qquad
        b_0=\lceil\log_2m\rceil,\qquad F=400.
\]
Choose $a_0$ left/right partitions of $\Lambda$ with the following
property: whenever $\alpha\ne\lambda$ and
$\{\alpha,\lambda\}\ne\{\gamma,\nu\}$, some partition places
$\alpha$ on the left, $\lambda$ on the right, and $\gamma,\nu$ on
the same side. This includes $\gamma=\nu$.

Assign each tag independently to either side, with equal probabilities.
A constraint is feasible even when some displayed tags coincide: the
excluded set equality ensures that putting $\gamma$ and $\nu$ together
does not force $\alpha$ and $\lambda$ together. At most four tag choices
therefore suffice to satisfy it, giving probability at least $1/16$.
There are no more than $m^4$ constraints. With $a_0$ independent
partitions, the expected number left unsatisfied is bounded by
\[
         m^4(15/16)^{a_0}
          \le m^4e^{-a_0/16}\le m^4(2m)^{-4}=1/16<1.
\]
Some choice therefore misses none. Give each partition one circle
coordinate, called a first coordinate.

Also choose $b_0$ partitions into Types~1 and~2 that separate every
two distinct tags, by assigning distinct binary strings of length
$b_0$ to the tags. Each such partition receives a batch of coordinates.
For every ordered list of $F$ directed pairs $(p,p')$ of distinct
roster slots whose underlying unordered edges are distinct and form
a forest, include a dedicated coordinate in this batch. The forest
restriction will allow arbitrary edge differences to be realized by
successive assignments of the roster values. Add spare coordinates if
necessary so that the batch has positive odd size.
There are finitely many lists, and a one-coordinate spare batch is
allowed if no list exists. These first coordinates and batches form
the finite coordinate set $C$ of $\mathcal T$. For $m=1$, there are
no Type partitions, but $a_0\ge1$ still supplies first coordinates.

\paragraph{First coordinates and their roster values.}
Enumerate $\mathcal P=\{p_0,\ldots,p_{r-1}\}$. In each first
coordinate choose $v_{p_j}=\theta2^j\pmod1$, where $\theta>0$ is
small enough that all the values lie in an arc of length less than
$1/8$. The ordered differences for distinct slots are nonzero and
pairwise distinct on the circle. Indeed a positive integer difference
$2^j-2^i=2^i(2^{j-i}-1)$ determines $i$ by its power of two, then
$j$ by its odd factor; the sign distinguishes the reverse orientation.
The small scale prevents collisions modulo one.

Since this is finite data, choose $0<\epsilon<1/100$ so small that
these nonzero differences are separated from one another and from
zero by more than $12\epsilon$, with a strict margin. If there are
no distinct-slot pairs, any sufficiently small positive $\epsilon$
works. In the first coordinate for a left/right partition, impose
the requirements
\[
\begin{array}{c|cc}
 \text{side of }\alpha & I_\alpha & J_\alpha\\ \hline
 \text{left}  & \|i\|_{\mathbb T}\le\epsilon
               & \|j\|_{\mathbb T}\ge4\epsilon\\
 \text{right} & \|i\|_{\mathbb T}\ge4\epsilon
               & \|j\|_{\mathbb T}\le\epsilon.
\end{array}
\]
An edge from a left tag to a right tag forces its total to have
distance at most $2\epsilon$ from zero. A same-side edge has one
near and one far summand, so its total has distance at least
$3\epsilon$ from zero. These possibilities are disjoint.

Every total $z$ with $\|z\|_{\mathbb T}>6\epsilon$ has a diagonal
decomposition in this coordinate, by assigning zero to the near
summand and $z$ to the far summand. Perturbing the first summand by
distance at most $\epsilon/2$, while perturbing the second by the
opposite amount, preserves both requirements. Thus the diagonal
decomposition has a uniform interior margin.

\paragraph{Type~1 batches.}
Let
\[
       C_I=\{0,1,3\}/6,\qquad C_J=\{2,4,5\}/6
                 \quad\text{in }\mathbb T.
\]
For a Type~1 tag, every coordinate of its $I$-vector must be within
$1/16$ of $C_I$, and every coordinate of its $J$-vector within $1/16$
of $C_J$. Their center sums contain all six multiples of $1/6$.
For an arbitrary total $z$ in one coordinate, choose a center sum
$c_I+c_J$ with a signed error $\delta$ of magnitude at most $1/12$.
The decomposition
\[
                   z=(c_I+\delta/2)+(c_J+\delta/2)
\]
puts each summand within $1/24$ of its center, leaving margin $1/48$
inside the allowed radius. It therefore remains valid when the first
summand is moved by at most $1/96$ and the second by the opposite
amount. These choices are independent in the batch coordinates.

On the other hand, the center sets have circular distance at least
$1/6$. Every difference of allowed Type~1 vectors satisfies
\begin{equation}\label{clock:type-one-difference}
              \|i_c-j_c\|_{\mathbb T}
                 \ge 1/6-2/16=1/24
                 \quad\text{in every batch coordinate }c.
\end{equation}

\paragraph{Type~2 batches.}
In a batch of $2s+1$ coordinates, a Type~2 tag requires, separately
for each of its $I$- and $J$-vectors, at least $s+1$ coordinates at
distance at most $1/100$ from zero. Suppose a total $z$ has one
coordinate $c_0$ at distance at most $1/200$ from zero. Set the first
summand to $z_{c_0}$ and the second to zero there. Partition the other
$2s$ coordinates into two groups of size $s$. In one group set the
first summand to zero, and in the other set the second summand to
zero; in each coordinate the remaining summand is determined by $z$.
Each vector has $s+1$ designated small coordinates. At every such
coordinate the margin to the threshold $1/100$ is at least $1/200$,
so moving the first vector by at most $1/400$ in the maximum metric
and the second oppositely preserves both majority conditions.

Conversely, the small-coordinate sets of two allowed Type~2 vectors
intersect. Their difference has at least one coordinate satisfying
\begin{equation}\label{clock:type-two-difference}
                            \|i_c-j_c\|_{\mathbb T}\le2/100.
\end{equation}
This includes $s=0$, a one-coordinate batch.

\paragraph{The graph and the interior condition.}
Define $I_\alpha$ and $J_\alpha$ by all their coordinate requirements.
If an off-diagonal edge $\alpha\to\lambda$ occurs, a separating
left/right partition excludes every competitor
$(\gamma,\nu)$ with $\{\gamma,\nu\}\ne\{\alpha,\lambda\}$.
Only its reversal could remain. If both orientations occurred, there
would be suitable elements of their four sets with
\[
 i_\alpha+j_\lambda=i_\lambda+j_\alpha,
 \qquad i_\alpha-j_\alpha=i_\lambda-j_\lambda.
\]
A Type partition separating $\alpha$ and $\lambda$ makes the common
difference satisfy both \eqref{clock:type-one-difference} and
\eqref{clock:type-two-difference}, contrary to $2/100<1/24$.
This proves~\textup{(i)}.

Let $W$ be the union of these obstructions for a total $z$:
\begin{enumerate}[label=\textup{(\alph*)}]
\item some first coordinate has distance at most $6\epsilon$ from zero;
\item some Type batch has no coordinate at distance at most $1/200$
from zero.
\end{enumerate}
If $z\notin W$, all the diagonal constructions above apply. They
combine across disjoint coordinate blocks. For example, the common
positive radius
\begin{equation}\label{clock:interior-radius}
                  \rho=\min\{\epsilon/2,1/96,1/400\}
\end{equation}
gives a ball in $I_\alpha\cap(z-J_\alpha)$ for every $\alpha$.
Indeed a perturbation of the first vector in this ball produces the
opposite perturbation of the second, and both remain allowed in every
block. This proves the interior assertion. Since $\Lambda$ is nonempty,
the graph alternative now makes the graph outside $W$ exactly the
full diagonal.

At this point the sets $I_\alpha,J_\alpha,W$ have all the required
compatibility and interior properties. The remaining choices concern the
roster alone. The first coordinates already have separated differences.
In a Type batch we use a different device: any sufficiently large set of
obstructed pairs would contain a forest, and one coordinate was reserved
specifically to prevent that forest from being obstructed all at once.

\paragraph{Roster values in dedicated batch coordinates.}
For the coordinate dedicated to the directed forest list
$((p_j,p'_j))_{j=0}^{F-1}$, prescribe
\begin{equation}\label{clock:forest-values}
                         v_{p'_j}-v_{p_j}=j/F\pmod1.
\end{equation}
Choose a root in each tree. Starting from any root value, traverse its
edges and assign the next vertex by the prescribed signed difference.
A forest has no cycle, so no vertex assignment can conflict with an
earlier one. Give arbitrary values to other vertices and to spare
coordinates. Since every coordinate makes its own assignments, these
choices together specify all points $v_p\in\mathcal T$.

For every coordinate translate $a_c$, one of the equally spaced
points $j/F$ is at distance at most $1/(2F)=1/800$ from $-a_c$.
If each roster point is moved by less than $\eta$ in the maximum
metric, each prescribed difference moves by less than $2\eta$.
Choose $\eta>0$ sufficiently small that
\begin{equation}\label{clock:perturbation-margin}
                     1/800+2\eta<1/200
\end{equation}
and so that all ordered distinct-slot differences in every first
coordinate remain separated from each other and zero by more than
$12\epsilon$. The latter is possible because their original finite
separations have a strict margin; a difference moves by at most
$2\eta$, and a separation between two differences by at most
$4\eta$. Thus, uniformly in every translate and every permitted
perturbation, each dedicated forest list contains a pair whose
translated difference is within $1/200$ of zero in its dedicated
coordinate.

\paragraph{Counting obstructions.}
Fix any permitted perturbation and any $a\in\mathcal T$. In a first
coordinate, at most one ordered distinct-slot pair can have its
translated difference within $6\epsilon$ of zero: two such
differences would be at distance at most $12\epsilon$ from each
other. Thus the first-coordinate obstructions cost at most $a_0$
pairs in total.

Fix one batch, and consider the directed pairs whose translated
differences have no coordinate within $1/200$ of zero in that batch.
If their underlying simple graph contains a forest with $F$ edges,
choose an available orientation of each edge and order the resulting
pairs. Its dedicated coordinate, by
\eqref{clock:perturbation-margin}, makes one of those pairs close to
zero, a contradiction. Consequently a spanning forest on the
nonisolated vertices has fewer than $F$ edges. If it has $f$ edges,
$v$ vertices, and $t$ components, then $f=v-t$ and $t\le v/2$;
hence $v\le2f<2F$. There are fewer than $4F^2$ directed pairs on
these vertices. In particular this batch contributes at most $4F^2$
pairs. The same argument covers a roster too small to have a dedicated
forest list.

Taking the union over coordinates and batches bounds the left-hand
side of \eqref{clock:roster-estimate} by $a_0+4F^2b_0$. Since
$\log2>1/2$,
\[
 a_0\le66\log(2m),\qquad
 b_0\le\frac{\log(2m)}{\log2}\le2\log(2m).
\]
It follows that
\[
 a_0+4F^2b_0
 \le(66+8\cdot400^2)\log(2m)
 <10^7\log(2m).
\]
This proves~\textup{(iii)} with one tolerance $\eta$ for all roster
points, all translates, and all coordinates. If the roster has one
element there are no distinct-slot pairs, so its estimates are
vacuous; if $m=1$ the Type-batch estimates are absent. All the other
constructions and the interior radius remain valid in these cases.
\end{proof}

To see how the clock will be applied, take a periodic additive map
$\varphi:\mathbb Z\to\mathcal T$. A prefix ending at $k$ can test
$\varphi(k-t)\in I_\alpha$, and a suffix starting there can independently
test $\varphi(q-k)\in J_\lambda$. Once $t,q,k$ refer to the same
positions, the pair lies in $\mathcal E_{\varphi(q-t)}$. The graph
alternative then gives either equality of tags or a unique unequal pair.
It says nothing by itself about whether the suffix chose the correct
anchor $q$; that is the purpose of the end guard in
Section~\ref{episode:section}. Finally, periodicity makes all clock
membership tests finite tables on integer residues. The torus is used
to construct the tables, not as an extra input alphabet.
